\section{Conjugación, orientación y retorno de las secciones materiales}
\label{vi:sec:conjugacion}

La identidad persistente de la sección~\ref{vi:sec:particula} admite
operaciones de inversión sobre niveles diferentes de su estructura.
La inversión celular actúa en el atlas; la conjugación dodecafásica
actúa en el registro de orientación; el levantamiento de un lazo actúa
en una fibra; una rotación espinorial actúa en un espacio de
representación. Especificar esos dominios permite componer las
operaciones y explicar qué información se conserva en cada retorno.

\subsection{Involución dodecafásica y paridad de los observables}
\label{vi:subsec:conjugacion-registro}

Sea $\mathcal T_{12}=\{-1,0,1\}^{12}$ el ambiente de registros
tríticos de doce posiciones. El dominio admisible del TPK es el
subconjunto con las compatibilidades de ruta ya construidas. La
operación sobre el ambiente es
\begin{equation}
 C_M:\mathcal T_{12}\longrightarrow\mathcal T_{12},\qquad
 (C_M\tau)_j=-\tau_{11-j},\quad 0\leq j<12.
 \label{vi:eq:cm}
\end{equation}
La restricción a registros admisibles utiliza la estabilidad de ese
dominio bajo la conjugación. El orden invertido y el cambio de signo
son componentes simultáneas de la operación.

\begin{proposition}[Paridad lineal y cuadrática del registro]
\label{vi:prop:paridad-registro}
Se cumple $C_M^2=I$. Si $a_j=a_{11-j}$ y
$L_a(\tau)=\sum_ja_j\tau_j$, entonces
$L_a(C_M\tau)=-L_a(\tau)$. Si una forma cuadrática
$Q_b(\tau)=\sum_{j,k}b_{jk}\tau_j\tau_k$ satisface
$b_{jk}=b_{11-j,11-k}$, entonces $Q_b(C_M\tau)=Q_b(\tau)$.
\end{proposition}
\begin{proof}
En cada posición,
$[C_M(C_M\tau)]_j=-(-\tau_{11-(11-j)})=\tau_j$.
Para el funcional lineal, sustituir $k=11-j$ en la suma da
$L_a(C_M\tau)=-\sum_ka_{11-k}\tau_k=-L_a(\tau)$.
En la expresión cuadrática los dos cambios de signo se cancelan y
la sustitución $(r,s)=(11-j,11-k)$ recupera $Q_b$ por la
simetría de sus coeficientes.
\end{proof}

La orientación lineal y las correlaciones de memoria tienen, por tanto,
leyes distintas. Incluso una composición de estos observables debe
conservar sus paridades: si $L_a$ es impar, su carácter exponencial
satisface
\begin{equation}
 \exp(L_a(C_M\tau))=\exp(L_a(\tau))^{-1}.
 \label{vi:eq:caracter-inverso}
\end{equation}
La igualdad expresa inversión multiplicativa. La invariancia de un
operador material completo se estudia sobre ese operador, con sus
canales y su representación, y se demuestra por la covariancia
correspondiente.

La inversión celular del teorema~\ref{vi:thm:selector-central} es
$\rho_c:\mathcal C_{81}\to\mathcal C_{81}$, mientras $C_M$ tiene
el dominio \eqref{vi:eq:cm}. Ambas son involuciones, pero para
compararlas se requiere el mapa de registro de la ruta y su ley de
transporte. Este dato evita sustituir la conjugación de una historia
por la inversión de una coordenada aislada del atlas.

\subsection{Levantamiento a antisecciones y conservación de la masa}
\label{vi:subsec:antisecciones}

La conjugación se aplica a la partícula mediante dos mapas
compatibles. El mapa de base $C_m^{\mathrm b}$ transporta el estado
$x_m$ a su estado conjugado $\bar x_m$; el levantamiento
$\widetilde C_m:\mathscr E_m(x)\to\mathscr E_m(\bar x)$
transporta sus secciones materiales. Sean $t_m$ el lector del
registro dodecafásico y $\bar p_{n,m}$ las restricciones en la
familia conjugada. Sus condiciones son conservación de la
admisibilidad, involutividad y compatibilidad:
\begin{equation}
 \begin{aligned}
 \rho_{n,m}C_n^{\mathrm b}&=C_m^{\mathrm b}\rho_{n,m},
 &t_m(C_m^{\mathrm b}x_m)&=C_Mt_m(x_m),\\
 \bar p_{n,m}\widetilde C_n&=\widetilde C_mp_{n,m},
 &\widetilde C_{m,\bar x}\widetilde C_{m,x}&=I.
 \end{aligned}
 \label{vi:eq:levantamiento-cm}
\end{equation}
En la primera fila, el lector publica el registro de orientación
de la historia. El transporte restante del estado acompaña esa
publicación. La segunda fila da las leyes del levantamiento;
los subíndices $x,\bar x$ explicitan sus dos fibras de partida.
Si $s$ es persistente, su antisección es
$\bar s=(\widetilde C_ms_m)_m$.

\begin{proposition}[Persistencia de la antisección]
\label{vi:prop:antiseccion-persistente}
Las condiciones \eqref{vi:eq:levantamiento-cm} transportan secciones
persistentes a secciones persistentes y restituyen la sección al
conjugar dos veces. Si además transportan las simetrías internas por
$g_m\mapsto\widetilde C_mg_m\widetilde C_m^{-1}$, inducen una
involución en las clases materiales de la
definición~\ref{vi:def:particula}.
\end{proposition}
\begin{proof}
La compatibilidad da
$\bar p_{n,m}\bar s_n=\widetilde C_mp_{n,m}s_n
=\widetilde C_ms_m=\bar s_m$.
La involutividad restituye $s$. Para $s'_m=g_ms_m$, sus
antisecciones se relacionan mediante
$\widetilde C_mg_m\widetilde C_m^{-1}$. La ley de restricción
en \eqref{vi:eq:levantamiento-cm} conserva la compatibilidad de esta
familia, que pertenece a los grupos internos por la condición del
enunciado. Así la construcción desciende a las clases.
\end{proof}

La representación de carga utiliza la parte impar del registro
material. Si $\mathfrak q$ es el carácter de carga declarado, su
ley es $\mathfrak q(\bar s)=-\mathfrak q(s)$. La carga elemental
que interviene al expresar esa coordenada en una recta dimensional es
la salida interna ya construida; la representación se aplica después
al estado, sin utilizar una carga objetivo para elegir su registro.

Para la masa se necesita la covariancia de su operador. Sean
$\widehat M_s$ y $\widehat M_{\bar s}$ los operadores autoadjuntos
de las dos realizaciones, y sea $\mathcal C$ su entrelazador
unitario o antiunitario. Su dominio se transporta mediante
$\mathcal C\operatorname{Dom}(\widehat M_s)
=\operatorname{Dom}(\widehat M_{\bar s})$. La condición operatoria es
\begin{equation}
 \mathcal C\widehat M_s\mathcal C^{-1}
 =\widehat M_{\bar s}.
 \label{vi:eq:covariancia-masa}
\end{equation}

\begin{proposition}[Conservación espectral bajo conjugación]
\label{vi:prop:masa-conjugada}
Bajo \eqref{vi:eq:covariancia-masa}, el transporte de una sección
propia de masa $m\in\mathbb R$ conserva ese autovalor. La medida
espectral se transporta mediante
$E_{\bar s}(B)=\mathcal C E_s(B)\mathcal C^{-1}$ para todo
boreliano real $B$. En particular, una realización escalar conserva
la masa al pasar a la antisección.
\end{proposition}
\begin{proof}
Si $\widehat M_s\psi=m\psi$, la covariancia da
$\widehat M_{\bar s}\mathcal C\psi
=\mathcal C(m\psi)=m\mathcal C\psi$; la última igualdad
vale también en el caso antiunitario porque $m$ es real. El cálculo
funcional real de los operadores autoadjuntos transporta sus
proyectores espectrales y prueba la segunda afirmación.
\end{proof}

En un registro de banda puede verse la misma condición en forma
elemental. Dada una lectura real $E_0(\tau)$, sean
$E_\pm=(E_0\pm E_0\circ C_M)/2$ y sea $\xi\in\{+1,-1\}$
la orientación de hoja. Entonces
\begin{equation}
 E_{\mathrm{band}}(\tau,\xi)=E_+(\tau)+\xi E_-(\tau),\qquad
 E_{\mathrm{band}}(C_M\tau,-\xi)=E_{\mathrm{band}}(\tau,\xi).
 \label{vi:eq:banda-paridad}
\end{equation}
La igualdad se obtiene usando que $E_+$ es par y $E_-$ impar.
Explica cómo puede conservarse la energía de banda al transportar
simultáneamente registro y hoja; su identificación con la masa de una
realización concreta utiliza \eqref{vi:eq:covariancia-masa}.

El registro electrónico
$\tau_e=(+,+,+,-,-,-,+,+,+,-,-,-)$ satisface $C_M\tau_e=\tau_e$.
Esta igualdad concierne a la palabra de doce posiciones. El par
$(\tau_e,\xi)$ conserva una orientación adicional que la
conjugación cambia. La distinción permite mantener el centro
electrónico único y, a la vez, las realizaciones conjugadas de carga;
una palabra fija no se convierte en una prueba de neutralidad.

\subsection{Holonomía de signo y banda de persistencia}
\label{vi:subsec:mobius}

Sea $\ell$ un lazo de retorno del lector material de una extensión
superviviente. Su grupo de vueltas es $\langle\ell\rangle\simeq
\mathbb Z$. Sobre este grupo se declara el carácter de signo
\begin{equation}
 \varepsilon:\langle\ell\rangle\longrightarrow\{+1,-1\},
 \qquad \varepsilon(\ell)=-1.
 \label{vi:eq:holonomia-signo}
\end{equation}
El grupo de vueltas registra cuántas veces se recorre el lazo; su
tipo es distinto de la fase finita $C_9$. El fibrado de intervalos
asociado al carácter es
\begin{equation}
 \mathcal M_\ell=([0,1]\times[-1,1])/
                  \bigl((0,u)\sim(1,-u)\bigr).
 \label{vi:eq:banda-mobius}
\end{equation}

\begin{theorem}[Retorno orientado sobre la banda]
\label{vi:thm:retorno-mobius}
El fibrado \eqref{vi:eq:banda-mobius} es una banda de Möbius.
Su transporte sobre $\ell$ cambia el signo de la fibra y sobre
$\ell^2$ lo restituye. Por tanto, un vector no nulo de esa línea
requiere dos vueltas para su retorno orientado. Esto conserva la
distinción entre retorno de orientación y retorno del estado completo.
\end{theorem}
\begin{proof}
La relación de extremos de \eqref{vi:eq:banda-mobius} es el mapa
lineal $u\mapsto-u$, con determinante negativo. Da la línea real
no orientable sobre el círculo y su correspondiente banda de
intervalos. Componer dos transportes produce $u\mapsto u$.
La memoria de la extensión no forma parte de ese cociente de signo;
su actualización continúa obedeciendo el transporte del TPK y
\eqref{vi:eq:persistencia-fase-memoria}.
\end{proof}

Una parametrización de una vuelta de la base mediante $2\pi$
representa la restitución de la cubierta mediante $4\pi$. Si el
lector temporal asigna $108$ pasos a la vuelta relevante, el
levantamiento orientado cierra a $216$ pasos. Estas asignaciones
requieren el lazo y el lector indicados; no dividen la historia en
dos partículas sucesivas de $108$ pasos.

\begin{proposition}[Descomposición de observables en la cubierta doble]
\label{vi:prop:observables-paridad}
Sea $p:\widetilde X\to X$ una cubierta principal doble, con
involución $\iota$. Para $f\in C^0(\widetilde X,\mathbb R)$,
\[
 f^+=\frac{f+f\circ\iota}{2},\qquad
 f^-=\frac{f-f\circ\iota}{2}
\]
dan $f=f^++f^-$, con $f^+\circ\iota=f^+$ y
$f^-\circ\iota=-f^-$. La primera parte desciende a $X$; la
segunda representa una sección de la línea de signo asociada. En
la línea de Möbius, toda sección global continua tiene algún cero.
\end{proposition}
\begin{proof}
Las identidades se obtienen sustituyendo $\iota^2=I$. La parte
invariante es constante en cada fibra y define una función en el
cociente. La parte antiinvariante satisface la ley de transición
de la representación de signo. En el círculo, una sección de esta
línea se representa por una función real continua $u$ en $[0,1]$
con $u(1)=-u(0)$. Si el extremo vale cero ya hay un cero; en caso
contrario, los extremos tienen signos opuestos y el teorema del
valor intermedio proporciona un cero interior.
\end{proof}

La descomposición distingue dos clases de observables sobre una
misma cubierta. Permite representar información impar de orientación
y datos pares conservados. Para aplicarla a un registro material se
utiliza el levantamiento \eqref{vi:eq:levantamiento-cm}; el carácter
de signo describe la línea transportada, no todos los datos de la
partícula.

\subsection{Retículo bidireccional y conjugación de los canales}
\label{vi:subsec:reticulo-bidireccional}

El registro angular entero publica dos coordenadas
\begin{equation}
 W_+=6n_A+s_C,\qquad W_-=6n_A-s_C,
 \qquad (n_A,s_C)\in\mathbb Z^2.
 \label{vi:eq:reticulo-ida-retorno}
\end{equation}
Los coeficientes $n_A,s_C$ son lecturas internas de ruta. El factor
$6$ pertenece a la normalización dodecafásica de ese registro.

\begin{proposition}[Índice del retículo e intercambio de las lecturas]
\label{vi:prop:indice-bidireccional}
La imagen de \eqref{vi:eq:reticulo-ida-retorno} tiene índice $12$
en $\mathbb Z^2$ y forma normal de Smith $\operatorname{diag}(1,12)$.
Se caracteriza por $W_++W_-\equiv0\pmod{12}$. La operación
$s_C\mapsto-s_C$ intercambia $W_+$ y $W_-$.
\end{proposition}
\begin{proof}
La matriz del mapa es
$\bigl(\begin{smallmatrix}6&1\\6&-1\end{smallmatrix}\bigr)$.
El máximo divisor común de sus entradas es $1$ y su determinante
es $-12$, lo que da los factores invariantes indicados. Las
coordenadas inversas son
$n_A=(W_++W_-)/12$ y $s_C=(W_+-W_-)/2$.
Si la suma es divisible por $12$, los dos enteros tienen la misma
paridad y ambas expresiones son enteras. La condición es también
necesaria por la suma de \eqref{vi:eq:reticulo-ida-retorno}.
La afirmación sobre el intercambio se verifica sustituyendo $-s_C$.
\end{proof}

En la representación angular real, las coordenadas ya generadas
$a=A_{\mathrm{rad}}$ y $b=C^*_{\mathrm{rad}}$ producen los
canales $q_\pm=\exp[-(a\pm b)]$. En el dominio $a>|b|$ ambos
pertenecen a $(0,1)$. Sea $R^*=R$, $R^2=I$ la graduación de
hojas y sea $C$ un intercambio unitario con $CRC^{-1}=-R$.
Para $T(b)=\exp[-(aI+bR)]$ se obtiene
\begin{equation}
 CT(b)C^{-1}=T(-b),\qquad
 CP_\pm C^{-1}=P_\mp,\quad P_\pm=(I\pm R)/2.
 \label{vi:eq:conjugacion-canales}
\end{equation}
En efecto, la conjugación transforma el exponente y cada término
de su serie convergente. Esta prueba transmite el intercambio a
los canales. El retículo entero y el operador real tienen dominios
distintos y publican la misma operación de intercambio a través de
sus mapas de lectura declarados.

\subsection{Representación espinorial del bloque electrónico}
\label{vi:subsec:spin-electron}

La construcción APP tridimensional de los proyectores de suma y
producto selecciona un plano principal electrónico por el modo
trítico $(1,-1,0)^3$. Su derivación y la evaluación de los
proyectores se incluyen entre los antecedentes del presente
artículo. En una base ortonormal de ese plano, sus restricciones
son
\begin{equation}
 P=\begin{pmatrix}1&0\\0&0\end{pmatrix},\qquad
 Q=\begin{pmatrix}1/4&\sqrt3/4\\\sqrt3/4&3/4\end{pmatrix},
 \qquad S=2P-I,\quad J=\frac4{\sqrt3}[P,Q].
 \label{vi:eq:bloque-electronico}
\end{equation}
Esta representación recibe los dos proyectores construidos y
conserva su orientación. La multiplicación da
$S^*=S$, $J^*=-J$, $S^2=I$, $J^2=-I$ y $SJ=-JS$.
En concreto, $S=\operatorname{diag}(1,-1)$ y
$J=\bigl(\begin{smallmatrix}0&1\\-1&0\end{smallmatrix}\bigr)$.

Sea $\mathscr S_e=E_e\otimes_{\mathbb R}\mathbb C$ la
complejificación del plano real. La unidad escalar $i$ pertenece
al cuerpo de esta representación; $J$ sigue siendo el
endomorfismo real procedente de los proyectores.

\begin{proposition}[Terna hermítica y representación de rotaciones]
\label{vi:prop:spin}
En $\mathscr S_e$, defínanse
\begin{equation}
 \sigma_1=SJ,\qquad \sigma_2=-iJ,\qquad \sigma_3=S.
 \label{vi:eq:terna}
\end{equation}
Estos operadores son hermíticos, de traza nula, y satisfacen
\begin{equation}
 \sigma_a\sigma_b=\delta_{ab}I+
 i\sum_c\epsilon_{abc}\sigma_c,
 \qquad L_a=\sigma_a/2,\qquad
 [L_a,L_b]=i\sum_c\epsilon_{abc}L_c,
 \quad\sum_aL_a^2=\tfrac34I.
 \label{vi:eq:spin-algebra}
\end{equation}
Constituyen la representación irreducible de espín $1/2$ sobre
el bloque electrónico.
\end{proposition}
\begin{proof}
Las relaciones de $S,J$ dan $(SJ)^*=SJ$, $(-iJ)^*=-iJ$ y
$\sigma_a^2=I$. Los productos cíclicos son
$(SJ)(-iJ)=iS$, $(-iJ)S=iSJ$ y $S(SJ)=J=i(-iJ)$;
invertir el orden cambia el signo. Esto prueba la identidad de
producto y los conmutadores. En la base de
\eqref{vi:eq:bloque-electronico} las tres matrices son
\[
 \begin{pmatrix}0&1\\1&0\end{pmatrix},\quad
 \begin{pmatrix}0&-i\\i&0\end{pmatrix},\quad
 \begin{pmatrix}1&0\\0&-1\end{pmatrix}.
\]
Su traza es cero y, junto con $I$, generan $M_2(\mathbb C)$.
Un subespacio invariante por las tres es invariante por toda esa
álgebra y sólo puede ser $0$ o $\mathscr S_e$. Finalmente,
sumar $L_a^2=I/4$ da el valor $3I/4$ del operador cuadrático.
\end{proof}

El reconocimiento de esta terna como matrices de Pauli es posterior
a \eqref{vi:eq:bloque-electronico}. La norma y el producto
$\langle X,Y\rangle=\tfrac12\operatorname{tr}(XY)$ en el
espacio real de matrices hermíticas de traza nula dan su espacio
de direcciones. Para $n\in\mathbb R^3$, $|n|=1$, sean
\begin{equation}
 H_e(n)=\sum_an_a\sigma_a,\qquad
 h_e(n)=\tfrac12H_e(n),\qquad
 \Pi_\pm(n)=\tfrac12(I\pm H_e(n)).
 \label{vi:eq:helicidad}
\end{equation}

\begin{proposition}[Descomposición direccional y doble retorno]
\label{vi:prop:helicidad}
Los operadores $\Pi_\pm(n)$ son proyectores ortogonales de
rango uno y $h_e(n)\Pi_\pm(n)=\pm\tfrac12\Pi_\pm(n)$.
El grupo de rotaciones alrededor de $n$ es
\begin{equation}
 U_n(\theta)=e^{-i\theta H_e(n)/2}
 =\cos(\theta/2)I-i\sin(\theta/2)H_e(n),\qquad
 U_n(2\pi)=-I,\quad U_n(4\pi)=I.
 \label{vi:eq:retorno-spin}
\end{equation}
\end{proposition}
\begin{proof}
La parte antisimétrica de \eqref{vi:eq:spin-algebra} se anula
al contraerse con $n_an_b$, por lo que $H_e(n)^2=I$.
De aquí se siguen la idempotencia, la ortogonalidad y la suma
$\Pi_++\Pi_-=I$. Son hermíticos y de traza uno, luego su
rango es uno. Multiplicar por $H_e(n)/2$ da los autovalores.
Separar las potencias pares e impares de la exponencial produce
\eqref{vi:eq:retorno-spin} y su ley de grupo. La normalización
angular se comprueba por
\[
 U_n(\theta)H_e(v)U_n(\theta)^*
 =H_e\bigl(v\cos\theta+(n\times v)\sin\theta
       +n(n\cdot v)(1-\cos\theta)\bigr),
\]
que resulta de multiplicar con la identidad de producto de la
terna. La acción inducida en las direcciones es una rotación de
ángulo $\theta$; el levantamiento espinorial lleva su mitad.
\end{proof}

Si $n$ se realiza como dirección de propagación, $h_e(n)$ es
la helicidad adimensional. La inversión de esa dirección cumple
$H_e(-n)=-H_e(n)$ y $\Pi_\pm(-n)=\Pi_\mp(n)$.
Esta operación cambia las etiquetas direccionales; la conjugación
de carga transporta el registro material mediante
\eqref{vi:eq:levantamiento-cm}. La quiralidad se define mediante
la graduación relativista que corresponda a la representación y
sus proyectores, distintos de \eqref{vi:eq:helicidad}.

La evaluación escalar electrónica actúa como
$\varepsilon_e^\beta I$ sobre $\mathscr S_e$ y conmuta con
la terna y con $\Pi_\pm(n)$. La información rotacional conserva
la estructura de esa evaluación sin introducir otro factor en
la masa. Por su parte, la banda de Möbius transporta una línea
real alrededor de un lazo y la holonomía nonádica transporta
fase y memoria. Los tres retornos pueden intervenir en una
misma realización; sus dominios y entrelazadores determinan la
composición precisa.
